\ExerciseChapter{主成分分析}

\label{chapter:1}

\section{往年题}

\par\addvspace{6pt}\noindent\begin{minipage}{\linewidth}

\subsection{单项选择题}

\Question{M15-1}{主成分分析与因子分析的关系}

关于主成分分析与因子分析的关系，正确的是：


\begin{choices}

\item 主成分分析法是提取公因子的常用方法之一。

\item 主成分分析和因子分析的解均不是唯一的。

\item 一般而言，主成分分析有利于变量关系解释，而因子分析有利于变量信息综合。

\item 主成分分析和因子分析均含参数估计和假设检验两部分。

\item 两者均要求数据来自正态分布，且变量间互不相关。

\end{choices}

\end{minipage}\par\vspace{4pt}

\par\addvspace{6pt}\noindent\begin{minipage}{\linewidth}

\Question{M17-3}{主成分与因子分析的异同}

关于主成分分析与因子分析的关系，按课程惯用口径，错误的是：


\begin{choices}

\item 主成分法是提取公因子的常用方法之一。

\item 主成分分析和因子分析的解都不是唯一的。

\item 通常主成分分析有利于变量信息综合，因子分析有利于变量关系解释。

\item 因子旋转可以改变单个共同因子的解释方差百分比。

\item 通常PCA用于描述；在适当模型条件下，可以对共同因子模型作假设检验。

\end{choices}

\end{minipage}\par\vspace{4pt}

\par\addvspace{6pt}\noindent\begin{minipage}{\linewidth}

\Question{M21-3}{主成分与因子分析的异同}

关于主成分分析与因子分析的关系，按课程惯用口径，错误的是：


\begin{choices}

\item 主成分法是提取公因子的常用方法之一。

\item 主成分分析和因子分析的解都不是唯一的。

\item 通常主成分分析有利于变量信息综合，因子分析有利于变量关系解释。

\item 因子旋转可以改变单个共同因子的解释方差百分比。

\item 通常PCA用于描述；在适当模型条件下，可以对共同因子模型作假设检验。

\end{choices}

\end{minipage}\par\vspace{4pt}

\par\addvspace{6pt}\noindent\begin{minipage}{\linewidth}

\Question{G17-4}{异常科目表现与方差解释率}

\begin{center}\small\setlength{\tabcolsep}{4pt}\renewcommand{\arraystretch}{1.18}\begin{tabular}{@{}lrrr@{}}
\toprule
姓名 & 语文 & 数学 & 英语 \\
\midrule

马云 & 84 & 88 & 80 \\
方舟子 & 95 & 99 & 95 \\
三毛 & 98 & 20 & 90 \\
王宝强 & 60 & 45 & 55 \\
赵薇 & 85 & 50 & 80 \\
\bottomrule\end{tabular}\end{center}

按特征值大于1保留单成分。若希望删除一条记录，使第一成分的方差解释比例更大，应删去谁？


\begin{choices}

\item 马云。

\item 方舟子。

\item 三毛。

\item 王宝强。

\item 赵薇。

\end{choices}

\end{minipage}\par\vspace{4pt}

\par\addvspace{6pt}\noindent\begin{minipage}{\linewidth}

\Question{G17-5}{综合评价排序的稳定性}

\begin{center}\small\setlength{\tabcolsep}{4pt}\renewcommand{\arraystretch}{1.18}\begin{tabular}{@{}lrrr@{}}
\toprule
姓名 & 语文 & 数学 & 英语 \\
\midrule

马云 & 84 & 88 & 80 \\
方舟子 & 95 & 99 & 95 \\
三毛 & 98 & 20 & 90 \\
王宝强 & 60 & 45 & 55 \\
赵薇 & 85 & 50 & 80 \\
\bottomrule\end{tabular}\end{center}

关于使用题中主成分、两个旋转因子及总分进行评价，下列说法正确的是：


\begin{choices}

\item 无论选用何种指标，方舟子永远第一。

\item 无论选用何种指标，王宝强永远倒数第一。

\item 无论选用何种指标，马云总比赵薇强。

\item 无论选用何种指标，三毛总比王宝强强。

\item 无论选用何种指标，赵薇总比王宝强强。

\end{choices}

\end{minipage}\par\vspace{4pt}

\par\addvspace{6pt}\noindent\begin{minipage}{\linewidth}

\Question{N19-6}{主成分的方差}

若以协方差阵的单位特征向量 \(a_k\) 定义 \(C_k=a_k^T(X-\mu)\)，则 \(\operatorname{Var}(C_k)\) 等于：


\begin{choices}

\item 对应特征值 \(\lambda_k\)。

\item 总等于1。

\item 总等于0。

\item 所有特征值的乘积。

\item 样本量。

\end{choices}

\end{minipage}\par\vspace{4pt}

\par\addvspace{6pt}\noindent\begin{minipage}{\linewidth}

\Question{N21-9}{相关阵PCA的总方差}

对 \(p\) 个非恒定变量作标准化，再对相关阵PCA，全部特征值之和为：


\begin{choices}

\item 1。

\item \(p\)。

\item \(p^2\)。

\item 样本数。

\item 0。

\end{choices}

\end{minipage}\par\vspace{4pt}

\par\addvspace{6pt}\noindent\begin{minipage}{\linewidth}

\Question{N21-10}{量纲改变与PCA}

仅把一个变量的单位从米改为厘米，且其他变量不变。关于PCA，通常正确的是：


\begin{choices}

\item 协方差阵PCA和相关阵PCA都一定改变。

\item 两者都一定不变。

\item 相关阵PCA不受此正比例单位换算影响，协方差阵PCA一般会改变。

\item 协方差阵PCA不变，相关阵一定改变。

\item 所有特征值都必变为100倍。

\end{choices}

\end{minipage}\par\vspace{4pt}

\par\addvspace{6pt}\noindent\begin{minipage}{\linewidth}

\Question{N21-11}{完整变换与降维}

保留全部单位正交主成分，对中心化样本点作变换后，下列哪项保持？


\begin{choices}

\item 欧氏距离。

\item 每个原坐标的数值。

\item 每个坐标的方差都相等。

\item 各主成分均值为原变量各自均值。

\item 所有样本点变成同一点。

\end{choices}

\end{minipage}\par\vspace{4pt}

\par\addvspace{6pt}\noindent\begin{minipage}{\linewidth}

\subsection{简答题}

\Question{PS20}{主成分变换与样本点的几何关系}

主成分分析改变了原样本空间中样本点的关系吗？Z变换会改变原样本空间中样本点的关系吗？为什么？


\worklines{3}

\end{minipage}\par\vspace{4pt}

\par\addvspace{6pt}\noindent\begin{minipage}{\linewidth}

\subsection{计算与分析题}

\Question{P20}{三科成绩：协方差阵与相关阵}

就下列成绩单，利用主成分分析计算加权平均分。共20分。

\begin{center}\small\setlength{\tabcolsep}{4pt}\renewcommand{\arraystretch}{1.18}\begin{tabular}{@{}lrrr@{}}
\toprule
ID & 语文 & 数学 & 英语 \\
\midrule

1 & 84 & 88 & 80 \\
2 & 95 & 99 & 95 \\
3 & 98 & 30 & 90 \\
4 & 60 & 55 & 55 \\
5 & 85 & 80 & 80 \\
\bottomrule\end{tabular}\end{center}

\begin{enumerate}
\def\labelenumi{\arabic{enumi}.}
\tightlist
\item
  对3科成绩做简单统计描述。（2分）
\item
  基于协方差阵计算3科权重及5人的加权平均分。（6分）
\item
  基于相关阵计算3科权重及5人的加权平均分。（6分）
\item
  比较第3、4号学生在两种方法下的得分，说明两种综合评分的关键区别。（6分）
\end{enumerate}


\worklines{3}

\end{minipage}\par\vspace{4pt}

\par\addvspace{6pt}\noindent\begin{minipage}{\linewidth}

\Question{P21}{三科成绩与招生评价}

利用下表做主成分分析（权重保留3位小数）。

\begin{center}\small\setlength{\tabcolsep}{4pt}\renewcommand{\arraystretch}{1.18}\begin{tabular}{@{}lrrr@{}}
\toprule
ID & 语文 & 数学 & 英语 \\
\midrule

1 & 84 & 88 & 80 \\
2 & 95 & 99 & 95 \\
3 & 98 & 20 & 90 \\
4 & 60 & 45 & 55 \\
5 & 85 & 80 & 80 \\
\bottomrule\end{tabular}\end{center}

\begin{enumerate}
\def\labelenumi{\arabic{enumi}.}
\tightlist
\item
  基于协方差阵计算3科权重。
\item
  基于相关阵计算3科权重。
\item
  一位同学的分数是\ldots\ldots，计算其两种加权平均分。
\item
  若你是理科院校招生人员，更愿采用哪种权重作为招生标准？说明理由。
\end{enumerate}

原稿没有记下第3问的三科分数。


\worklines{3}

\end{minipage}\par\vspace{4pt}

\par\addvspace{6pt}\noindent\begin{minipage}{\linewidth}

\Question{P19}{父母与子女身高的主成分回归}

对父母与子女身高数据进行主成分分析。共20分。

\begin{enumerate}
\def\labelenumi{\arabic{enumi}.}
\tightlist
\item
  写出主成分与原变量的函数关系式。
\item
  由主成分回归方程求原变量尺度的回归方程。
\end{enumerate}

以父母身高为自变量、子女身高为结局。原题所指身高数据未随往年题附上。


\worklines{3}

\end{minipage}\par\vspace{4pt}

\Question{PU}{从SPSS输出计算权重与主成分得分}

下列为三科成绩的SPSS分析结果。采用相关阵主成分分析，提取3个成分，各变量提取共同度均为1。共20分。

\begin{center}\small\setlength{\tabcolsep}{4pt}\renewcommand{\arraystretch}{1.18}\begin{tabular}{@{}lrrr@{}}
\toprule
变量 & 均值 & 标准差 & N \\
\midrule

语文 & 79 & 12.4 & 5 \\
英语 & 80 & 7.9 & 5 \\
数学 & 67 & 18.2 & 5 \\
\bottomrule\end{tabular}\end{center}

\begin{center}\small\setlength{\tabcolsep}{4pt}\renewcommand{\arraystretch}{1.18}\begin{tabular}{@{}lrrr@{}}
\toprule
主成分 & 特征值 & 贡献率\% & 累计\% \\
\midrule

1 & 2.634 & 87.789 & 87.789 \\
2 & 0.308 & 10.272 & 98.061 \\
3 & 0.058 & 1.939 & 100 \\
\bottomrule\end{tabular}\end{center}

\begin{center}\small\setlength{\tabcolsep}{4pt}\renewcommand{\arraystretch}{1.18}\begin{tabular}{@{}lrrrrrr@{}}
\toprule
变量 & 载荷1 & 载荷2 & 载荷3 & 得分系数1 & 得分系数2 & 得分系数3 \\
\midrule

语文 & 0.9 & 0.43 & 0.07 & 0.34 & 1.4 & 1.25 \\
英语 & 0.93 & -0.34 & 0.13 & 0.35 & -1.11 & 2.23 \\
数学 & 0.98 & -0.07 & -0.19 & 0.37 & -0.23 & -3.26 \\
\bottomrule\end{tabular}\end{center}

\begin{enumerate}
\def\labelenumi{\arabic{enumi}.}
\tightlist
\item
  计算权重之和为1的三科统计学权重。
\item
  学生A语文62、英语79、数学91，计算其简单平均分和加权平均分。
\item
  计算学生A第一主成分得分。
\item
  计算第一主成分得分与所求加权平均分的相关系数。
\end{enumerate}


\worklines{3}

\par\addvspace{6pt}\noindent\begin{minipage}{\linewidth}

\Question{P17}{成绩综合评分及个体评价}

成绩如下，数据作Z变换后，以主成分法提取两个公因子，并作varimax旋转。

\begin{center}\small\setlength{\tabcolsep}{4pt}\renewcommand{\arraystretch}{1.18}\begin{tabular}{@{}lrrr@{}}
\toprule
姓名 & 语文 & 数学 & 英语 \\
\midrule

马云 & 84 & 88 & 80 \\
方舟子 & 95 & 99 & 95 \\
三毛 & 98 & 20 & 90 \\
王宝强 & 60 & 45 & 55 \\
赵薇 & 85 & 50 & 80 \\
\bottomrule\end{tabular}\end{center}

\begin{enumerate}
\def\labelenumi{\arabic{enumi}.}
\tightlist
\item
  用主成分分析获得3科加权平均分的权重，写出步骤。（Q6，10分）
\item
  A同学语文82、数学63、英语81，计算其加权平均分。（Q7，5分）
\item
  根据主成分及因子分析，对A的学业水平作简要评价。（Q8，5分）
\end{enumerate}


\worklines{3}

\end{minipage}\par\vspace{4pt}

\Question{P13}{空气污染综合指数}

下列为广州市某月3种空气污染物指数的主成分分析输出。共20分。

\begin{enumerate}
\def\labelenumi{\arabic{enumi}.}
\tightlist
\item
  分别从相关阵与协方差阵结果计算空气污染综合指数的统计学权重。
\item
  哪种结果更符合实际？说明理由。
\end{enumerate}

\begin{center}\small\setlength{\tabcolsep}{4pt}\renewcommand{\arraystretch}{1.18}\begin{tabular}{@{}lrrr@{}}
\toprule
变量 & 均值 & 标准差 & N \\
\midrule

TSP & 104.28 & 53.167 & 61 \\
SO2 & 41.62 & 21.503 & 61 \\
NOx & 183.26 & 68.898 & 61 \\
\bottomrule\end{tabular}\end{center}

\begin{center}\small\setlength{\tabcolsep}{4pt}\renewcommand{\arraystretch}{1.18}\begin{tabular}{@{}lrrrrrr@{}}
\toprule
成分 & 相关阵\(\lambda\) & 贡献\% & 累计\% & 协方差阵\(\lambda\) & 贡献\% & 累计\% \\
\midrule

1 & 2.169 & 72.306 & 72.306 & 6140.765 & 76.415 & 76.415 \\
2 & 0.606 & 20.215 & 92.521 & 1722.108 & 21.43 & 97.845 \\
3 & 0.224 & 7.479 & 100 & 173.167 & 2.155 & 100 \\
\bottomrule\end{tabular}\end{center}

相关阵：载荷与成分得分系数。

\begin{center}\small\setlength{\tabcolsep}{4pt}\renewcommand{\arraystretch}{1.18}\begin{tabular}{@{}lrrrrrr@{}}
\toprule
变量 & 载荷1 & 载荷2 & 载荷3 & 得分1 & 得分2 & 得分3 \\
\midrule

TSP & 0.741 & 0.671 & 0.031 & 0.341 & 1.107 & 0.139 \\
SO2 & 0.892 & -0.316 & 0.324 & 0.411 & -0.521 & 1.446 \\
NOx & 0.909 & -0.237 & -0.344 & 0.419 & -0.391 & -1.532 \\
\bottomrule\end{tabular}\end{center}

协方差阵：原尺度载荷与重标度载荷。

\begin{center}\small\setlength{\tabcolsep}{4pt}\renewcommand{\arraystretch}{1.18}\begin{tabular}{@{}lrrrrrr@{}}
\toprule
变量 & 原1 & 原2 & 原3 & 重标1 & 重标2 & 重标3 \\
\midrule

TSP & 39.208 & 35.906 & -0.419 & 0.737 & 0.675 & -0.008 \\
SO2 & 16.883 & -3.685 & 12.797 & 0.785 & -0.171 & 0.595 \\
NOx & 65.715 & -20.476 & -3.038 & 0.954 & -0.297 & -0.044 \\
\bottomrule\end{tabular}\end{center}

协方差阵输出的标准化成分得分系数：

\begin{center}\small\setlength{\tabcolsep}{4pt}\renewcommand{\arraystretch}{1.18}\begin{tabular}{@{}lrrr@{}}
\toprule
变量 & 1 & 2 & 3 \\
\midrule

TSP & 0.339 & 1.109 & -0.129 \\
SO2 & 0.059 & -0.046 & 1.589 \\
NOx & 0.737 & -0.819 & -1.209 \\
\bottomrule\end{tabular}\end{center}

重标度后提取平方载荷和为(2.070, 0.574, 0.356)，贡献率为(69.001\%,19.126\%,11.873\%)，累计为(69.001\%,88.127\%,100\%)。


\worklines{3}

\Question{P15}{用成绩因子分析结果评价新学生}

某人的成绩为语文70、数学60、英语90。根据下列分析输出，采用适当方法计算其加权平均分。

\begin{center}\small\setlength{\tabcolsep}{4pt}\renewcommand{\arraystretch}{1.18}\begin{tabular}{@{}lrrr@{}}
\toprule
变量 & 均值 & 标准差 & 共同度 \\
\midrule

语文 & 82.2 & 11.946 & 0.978 \\
数学 & 75.2 & 19.942 & 1 \\
英语 & 80.2 & 7.95 & 0.979 \\
\bottomrule\end{tabular}\end{center}

\begin{center}\small\setlength{\tabcolsep}{4pt}\renewcommand{\arraystretch}{1.18}\begin{tabular}{@{}lrrrrrr@{}}
\toprule
成分 & 初始\(\lambda\) & 贡献\% & 累计\% & 旋转\(\lambda\) & 贡献\% & 累计\% \\
\midrule

1 & 2.223 & 74.109 & 74.109 & 1.912 & 63.736 & 63.736 \\
2 & 0.734 & 24.462 & 98.57 & 1.045 & 34.835 & 98.57 \\
3 & 0.043 & 1.43 & 100 & & & \\
\bottomrule\end{tabular}\end{center}

\begin{center}\small\setlength{\tabcolsep}{4pt}\renewcommand{\arraystretch}{1.18}\begin{tabular}{@{}lrrrrrr@{}}
\toprule
变量 & 初始1 & 初始2 & 旋转1 & 旋转2 & 得分1 & 得分2 \\
\midrule

语文 & 0.963 & -0.225 & 0.959 & 0.24 & 0.525 & -0.075 \\
数学 & 0.632 & 0.775 & 0.208 & 0.978 & -0.229 & 1.069 \\
英语 & 0.947 & -0.289 & 0.974 & 0.176 & 0.558 & -0.155 \\
\bottomrule\end{tabular}\end{center}


\worklines{3}

\clearpage\section{补充练习}

本节为配合正文新编的补充练习。除注明沿用正文例题外，数值与研究摘要均为教学设定。题型与作答方式沿用本章往年题，解析见章末。

\par\addvspace{6pt}\noindent\begin{minipage}{\linewidth}

\subsection{单项选择题}

\Question{SUP-P1}{特征向量、相关系数与标准成分得分}

对3个标准化变量作相关阵主成分分析。第一特征值为 \(\lambda_1=2.25\)，相应单位特征向量中 \(Z_1\) 的系数为 \(a_{11}=0.50\)。记 \(C_1=a_1^TZ\)，\(F_1=C_1/\sqrt{\lambda_1}\)。下列哪组数值依次为 \(\operatorname{Corr}(Z_1,C_1)\) 和 \(F_1\) 中 \(Z_1\) 的系数？


\begin{choices}

\item \((0.50,0.50)\)。

\item \((0.75,1/3)\)。

\item \((1.125,0.50)\)。

\item \((1/3,0.75)\)。

\item \((0.75,0.75)\)。

\end{choices}

\end{minipage}\par\vspace{4pt}

\par\addvspace{6pt}\noindent\begin{minipage}{\linewidth}

\subsection{简答题}

\Question{SUP-P2}{Bartlett球形检验与KMO的联合判断}

某研究拟将6项数值指标合成为一个综合指标，样本量180。Bartlett球形检验 \(P<0.001\)，KMO为0.42。研究者称：``球形检验显著，说明非常适合作主成分分析，并且只需保留第一主成分。''请评价这段话，写出球形检验的零假设，并说明还需查看哪些正文所讲的信息。


\worklines{3}

\end{minipage}\par\vspace{4pt}

\par\addvspace{6pt}\noindent\begin{minipage}{\linewidth}

\subsection{计算与分析题}

\Question{SUP-P3}{两科成绩的主成分与百分位评价}

对语文 \(X_1\)、数学 \(X_2\) 作相关阵主成分分析。两科均值分别为70、80，标准差分别为10、5，相关矩阵及单位特征向量如下： \[R=\begin{pmatrix}1&0.8\\0.8&1\end{pmatrix},\qquad
a_1=\frac1{\sqrt2}\begin{pmatrix}1\\1\end{pmatrix},\quad
a_2=\frac1{\sqrt2}\begin{pmatrix}1\\-1\end{pmatrix}.\] 学生A的语文、数学成绩为90、85。共15分。

\begin{enumerate}
\def\labelenumi{\arabic{enumi}.}
\tightlist
\item
  写出两个主成分的表达式、方差及各自贡献率。（4分）
\item
  计算A的两个主成分得分，并解释两个得分分别反映什么。（4分）
\item
  若两科成绩联合服从二元正态分布，按正文方法将第一主成分得分换算为正态百分位。已知 \(\Phi(1.5811)=0.9431\)。（4分）
\item
  只保留第一成分时，丢失了哪一方面的信息？能否根据``两个成分不相关''直接断言任意分布下二者独立？（3分）
\end{enumerate}


\worklines{3}

\end{minipage}\par\vspace{4pt}

\par\addvspace{6pt}\noindent\begin{minipage}{\linewidth}

\Question{SUP-P4}{胎儿受精龄：还原主成分回归方程}

沿用正文的胎儿外型测量例题：\(Y\) 为受精龄（周），\(X_1\) 为身高（cm），\(X_2\) 为头围（cm），\(X_3\) 为体重（g）。已知样本量为22，摘要如下。

\begin{center}\small\setlength{\tabcolsep}{4pt}\renewcommand{\arraystretch}{1.18}\begin{tabular}{@{}lrr@{}}
\toprule
指标 & 均值 & 标准差 \\
\midrule

\(X_1\) & 33.05 & 9.71 \\
\(X_2\) & 23.26 & 6.86 \\
\(X_3\) & 936.9 & 690.3 \\
\bottomrule\end{tabular}\end{center}

令 \(Z_j=(X_j-\bar X_j)/s_j\)，正文给出的结果为 \[C_1=0.58Z_1+0.58Z_2+0.57Z_3,\] \[C_2=-0.42Z_1-0.39Z_2+0.82Z_3,\] \[\widehat Y=23.73+3.88C_1+3.10C_2.\] 本题按以上已显示的系数精度计算。共15分。

\begin{enumerate}
\def\labelenumi{\arabic{enumi}.}
\tightlist
\item
  将回归式还原为 \(\widehat Y=b_0+b_1X_1+b_2X_2+b_3X_3\)。（8分）
\item
  某胎儿 \((X_1,X_2,X_3)=(40,28,1500)\)，计算其预测受精龄。（4分）
\item
  为什么主成分回归可缓解自变量共线性？只保留部分成分是否保证预测误差一定更小？（3分）
\end{enumerate}


\worklines{3}

\end{minipage}\par\vspace{4pt}
\clearpage\section{习题解析}

\begin{keyidea}[核对方式]按题号核对。补写题目、原稿歧义及复算约定列在对应解析中。\end{keyidea}

\Answer{M15-1}{主成分分析与因子分析的关系}

\textbf{答案：A。}

按课程的``主成分法提取因子''口径选A。PCA通常用于综合与降维，共同因子模型偏重相关结构解释。须补充：PCA特征向量有整体正负号任意性，重根时基底也不唯一；因此B若按严格数学意义理解有歧义。教材所称PCA``唯一''通常已约定符号且特征根互异，不能忽略这些条件。


\Answer{M17-3}{主成分与因子分析的异同}

\textbf{答案：B（课程口径）。}

原卷以PCA在给定互异特征根、规范化和符号约定后确定，而因子模型仍有旋转不确定性为比较，预期选B。但不加这些约定时，PCA也有符号和重根基底的不唯一性，B并非严格错误。保留原题用于识别考试口径，同时明确数学限制。


\Answer{M21-3}{主成分与因子分析的异同}

\textbf{答案：B（课程口径）。}

原卷以PCA在给定互异特征根、规范化和符号约定后确定，而因子模型仍有旋转不确定性为比较，预期选B。但不加这些约定时，PCA也有符号和重根基底的不唯一性，B并非严格错误。保留原题用于识别考试口径，同时明确数学限制。


\Answer{G17-4}{异常科目表现与方差解释率}

\textbf{答案：C。}

分别删去每人后重新计算相关阵第一特征值与3的比值；三毛的数学与语文英语表现明显不一致，删除后综合相关结构更强。逐人删除结果已附。不能将此练习变成在实际研究中为了提高解释率任意删人。


\Answer{G17-5}{综合评价排序的稳定性}

\textbf{答案：E（限题列指标及正向约定）。}

按本题4种实际指标复算，方舟子第一主成分与数学因子最高，但第一语文英语因子最高的是三毛，所以A不成立；第二因子最差是三毛，B不成立；马云第一因子低于赵薇，C不成立；三毛第二因子低于王宝强，D不成立。E在本次正向因子约定及列举的4种指标中成立，可按命题意图选E；但``何种指标''\,``永远''若按无限范围理解都过强，改变方向或评价目标可改变排序。


\Answer{N19-6}{主成分的方差}

\textbf{答案：A。}

由特征方程，方差为 \(a_k^T\Sigma a_k=\lambda_k\)。


\textbf{补写说明。} 2019卷该部分仅有总题数和部分回忆，整题内容未记录。本题为按课程知识点新编的补充练习，不是恢复的往年原题；编号只用于本包覆盖核对。


\Answer{N21-9}{相关阵PCA的总方差}

\textbf{答案：B。}

特征值之和等于迹，相关阵对角线共p个1。


\textbf{补写说明。} 2021卷该部分仅有总题数和部分回忆，整题内容未记录。本题为按课程知识点新编的补充练习，不是恢复的往年原题；编号只用于本包覆盖核对。


\Answer{N21-10}{量纲改变与PCA}

\textbf{答案：C。}

标准化消除了该正比例尺度变化，协方差阵则改变相应方差与协方差。


\textbf{补写说明。} 2021卷该部分仅有总题数和部分回忆，整题内容未记录。本题为按课程知识点新编的补充练习，不是恢复的往年原题；编号只用于本包覆盖核对。


\Answer{N21-11}{完整变换与降维}

\textbf{答案：A。}

正交变换保距；降维后该结论一般不再成立。


\textbf{补写说明。} 2021卷该部分仅有总题数和部分回忆，整题内容未记录。本题为按课程知识点新编的补充练习，不是恢复的往年原题；编号只用于本包覆盖核对。


\Answer{PS20}{主成分变换与样本点的几何关系}

保留全部成分时，协方差阵PCA的中心化和正交旋转保留欧氏距离及夹角；若截取前几个成分，投影通常改变距离并丢弃信息。相关阵PCA先作Z标准化，各坐标按不同标准差缩放，一般会改变原空间的距离和角度；仅在各尺度相同等特殊条件下保留其相对形状。故须区分``完整正交变换''和``降维投影''。


\Answer{P20}{三科成绩：协方差阵与相关阵}

\textbf{权重约定。} 记最大特征值对应、整体正向的单位特征向量为 \(a_1\)。按课程参考代码将 \(w_j=a_{j1}/\sum_l a_{l1}\) 作为原分数的权重，计算 \(W=\sum_jw_jX_j\)。相关阵的主成分本身则是 \(C_1=\sum_ja_{j1}(X_j-\bar X_j)/s_j\)，两者并不相同。如果要求原量纲综合分与相关阵第一主成分严格线性等价，应将 \(a_{j1}/s_j\) 再归一化作为权重，并说明约定。

\textbf{统计描述。}

\begin{center}\small\setlength{\tabcolsep}{4pt}\renewcommand{\arraystretch}{1.18}\begin{tabular}{@{}lrrr@{}}
\toprule
统计量 & 语文 & 数学 & 英语 \\
\midrule

均值 & 84.400 & 70.400 & 80.000 \\
标准差 & 14.943 & 27.790 & 15.411 \\
\bottomrule\end{tabular}\end{center}

\textbf{权重（语文、数学、英语顺序）。} 协方差阵为(0.107726, 0.718583, 0.173691)，相关阵为(0.427731, 0.133767, 0.438502)。

\begin{center}\small\setlength{\tabcolsep}{4pt}\renewcommand{\arraystretch}{1.18}\begin{tabular}{@{}lrr@{}}
\toprule
ID & 协方差加权分 & 相关阵加权分 \\
\midrule

1 & 86.1796 & 82.7811 \\
2 & 97.8743 & 95.5351 \\
3 & 47.7468 & 85.3958 \\
4 & 55.5386 & 57.1387 \\
5 & 80.5386 & 82.1387 \\
\bottomrule\end{tabular}\end{center}

协方差阵中数学的方差较大，权重约71.9\%，第3人数学低分受到更强惩罚，得分低于第4人；相关阵消除了各科尺度差异，第3人的语文、英语优势得以体现。由协方差阵还是相关阵出发，应根据评价目的、变量量纲和变异的实际含义选择，不宜仅凭哪种排序更合意。


\Answer{P21}{三科成绩与招生评价}

\textbf{权重约定。} 记最大特征值对应、整体正向的单位特征向量为 \(a_1\)。按课程参考代码将 \(w_j=a_{j1}/\sum_l a_{l1}\) 作为原分数的权重，计算 \(W=\sum_jw_jX_j\)。相关阵的主成分本身则是 \(C_1=\sum_ja_{j1}(X_j-\bar X_j)/s_j\)，两者并不相同。如果要求原量纲综合分与相关阵第一主成分严格线性等价，应将 \(a_{j1}/s_j\) 再归一化作为权重，并说明约定。

协方差阵权重为 \textbf{(0.088, 0.767, 0.145)}；相关阵权重为 \textbf{(0.415, 0.159, 0.427)}。实际计算宜保留未舍入权重。

第3问设语文、数学、英语为 \((c,m,e)\)： \[W_S=0.08766886c+0.76723953m+0.14509161e,\] \[W_R=0.4146077c+0.1586438m+0.4267485e.\] 原分数缺失，无法还原该人的唯一数值结果。

第4问可按课程意图选协方差权重，因为它在本样本中更强调数学。但这是数学分数变异较大的结果，不是PCA自动识别了``理科能力''。招生权重还应经过专业论证、效度与稳定性检验。


\Answer{P19}{父母与子女身高的主成分回归}

若从父母身高的相关阵求得特征向量 \(a_{jk}\)，则 \[C_k=\sum_{j=1}^{2}a_{jk}\frac{X_j-\bar X_j}{s_j}.\] 保留 \(r\) 个主成分后的回归式为 \(\hat Y=\gamma_0+\sum_{k=1}^r\gamma_kC_k\)。代回可得 \[b_j=\frac{\sum_k\gamma_ka_{jk}}{s_j},\qquad b_0=\gamma_0-\sum_jb_j\bar X_j.\] 因此 \(\hat Y=b_0+b_1X_1+b_2X_2\)。若仅保留第一主成分，求和仅有 \(k=1\)。协方差阵主成分只作中心化时，不除以 \(s_j\)。原数据缺失，不能把``身长、头围、体重、孕周''的另一份数据当作本题身高数据。源码中附通用转换步骤；须用本题数据估计数值系数。


\Answer{PU}{从SPSS输出计算权重与主成分得分}

\textbf{权重约定。} 记最大特征值对应、整体正向的单位特征向量为 \(a_1\)。按课程参考代码将 \(w_j=a_{j1}/\sum_l a_{l1}\) 作为原分数的权重，计算 \(W=\sum_jw_jX_j\)。相关阵的主成分本身则是 \(C_1=\sum_ja_{j1}(X_j-\bar X_j)/s_j\)，两者并不相同。如果要求原量纲综合分与相关阵第一主成分严格线性等价，应将 \(a_{j1}/s_j\) 再归一化作为权重，并说明约定。

按给出的第一列得分系数归一化，\(w=(0.34,\allowbreak 0.35,\allowbreak 0.37)/1.06=(0.320755,\allowbreak 0.330189,\allowbreak 0.349057)\)；粗略写(0.32, 0.33, 0.35)亦可，但结果须与精度一致。简单均分 \(232/3=77.333\)；按未再舍入权重得加权分约77.736。

SPSS得分系数列对应方差约为1的标准成分得分： \[F_1=.34\frac{62-79}{12.4}+.35\frac{79-80}{7.9}+.37\frac{91-67}{18.2}\approx-0.02252.\] 若采用单位特征向量定义的、方差为 \(\lambda_1\) 的主成分，则 \(C_1=\sqrt{2.634}F_1\approx-0.03655\)。需说明得分采用哪种归一化。

\textbf{第4问不能直接回答1。} 将相关阵特征向量归一化后乘原始成绩，得到的原分加权分一般不与标准化主成分完全相关。设原始协方差为 \(S=D_sRD_s\)，\(b=D_s^{-1}a_1\)，则 \[\operatorname{Corr}(C_1,W)=\frac{b^TSw}{\sqrt{(b^TSb)(w^TSw)}}.\] 只有权重与 \(D_s^{-1}a_1\) 成比例等相容条件下才为1。原输出数值经过舍入，精确相关系数需完整协方差阵或原数据；也可明确改用与主成分等价的权重，再得到相关系数1。


\Answer{P17}{成绩综合评分及个体评价}

\textbf{权重约定。} 记最大特征值对应、整体正向的单位特征向量为 \(a_1\)。按课程参考代码将 \(w_j=a_{j1}/\sum_l a_{l1}\) 作为原分数的权重，计算 \(W=\sum_jw_jX_j\)。相关阵的主成分本身则是 \(C_1=\sum_ja_{j1}(X_j-\bar X_j)/s_j\)，两者并不相同。如果要求原量纲综合分与相关阵第一主成分严格线性等价，应将 \(a_{j1}/s_j\) 再归一化作为权重，并说明约定。

相关阵第一特征值为2.055585，贡献率约68.519\%；第一特征向量取正向后为(0.673130, 0.255288, 0.694063)，课程权重为 \textbf{(0.414877, 0.157344, 0.427779)}。A同学加权分约 \textbf{78.583}。

A的语文略低于样本均值84.4，数学高于均值60.4，英语高于均值80；旋转后可用``语文英语''和``数学''两维描述，因子得分见配套计算结果。平均分、主成分、旋转因子衡量的方向不同；不要把仅5人的样本排序当稳定的能力评价。


\Answer{P13}{空气污染综合指数}

\textbf{权重约定。} 记最大特征值对应、整体正向的单位特征向量为 \(a_1\)。按课程参考代码将 \(w_j=a_{j1}/\sum_l a_{l1}\) 作为原分数的权重，计算 \(W=\sum_jw_jX_j\)。相关阵的主成分本身则是 \(C_1=\sum_ja_{j1}(X_j-\bar X_j)/s_j\)，两者并不相同。如果要求原量纲综合分与相关阵第一主成分严格线性等价，应将 \(a_{j1}/s_j\) 再归一化作为权重，并说明约定。

相关阵第一载荷列具有共同的 \(\sqrt{\lambda_1}\) 因子，可用 \((.741, .892, .909)\) 归一化，约得 \textbf{(0.2915, 0.3509, 0.3576)}。协方差阵用原尺度载荷 \((39.208, 16.883, 65.715)\) 归一化，约得 \textbf{(0.3219, 0.1386, 0.5395)}。所附61日数据可复算更高精度结果。

不能直接归一化协方差阵输出的``标准化得分系数''(.339, .059, .737)，那是另一尺度。若希望消除不同污染物指数的波动尺度差异，相关阵较合适；若指数已按可比风险标准定义，且绝对变异有实际意义，可论证采用协方差阵。实际环境风险权重还涉及危害与政策目的，PCA仅提供统计权重。


\Answer{P15}{用成绩因子分析结果评价新学生}

\textbf{权重约定。} 记最大特征值对应、整体正向的单位特征向量为 \(a_1\)。按课程参考代码将 \(w_j=a_{j1}/\sum_l a_{l1}\) 作为原分数的权重，计算 \(W=\sum_jw_jX_j\)。相关阵的主成分本身则是 \(C_1=\sum_ja_{j1}(X_j-\bar X_j)/s_j\)，两者并不相同。如果要求原量纲综合分与相关阵第一主成分严格线性等价，应将 \(a_{j1}/s_j\) 再归一化作为权重，并说明约定。

若取相关阵第一主成分的统计权重，可归一化第一初始载荷(.963, .632, .947)，得权重约(0.378835, 0.248623, 0.372541)，其原分加权平均约 \textbf{74.965}。

另一种是因子综合得分：先根据本题均值和标准差求 \(Z\)，再由得分系数得到 \(\widehat F_1\approx0.3262\)、\(\widehat F_2\approx-0.9293\)。按旋转后方差贡献比例加权， \[F_{\rm comp}=\frac{1.912}{1.912+1.045}\widehat F_1+\frac{1.045}{1.912+1.045}\widehat F_2\approx-0.1175.\] 这个结果是标准化综合得分，不是百分制平均分；若需转换回百分制，需要另行规定基准和尺度。题目未唯一规定加权规则，应明确所用方法。

\subsection{补充练习解析}

\Answer{SUP-P1}{特征向量、相关系数与标准成分得分}

\textbf{答案：B。}

相关阵PCA中，\(\operatorname{Corr}(Z_1,C_1)=a_{11}\sqrt{\lambda_1}=0.50\times1.50=0.75\)。把 \(C_1\) 除以其标准差后，\(Z_1\) 的系数为 \(a_{11}/\sqrt{\lambda_1}=1/3\)。三种系数所对应的量不同，读取输出时应先核对表名与得分的标准化方式。




\Answer{SUP-P2}{Bartlett球形检验与KMO的联合判断}

球形检验的零假设是总体相关矩阵为单位阵 \(H_0:R=I_6\)。显著结果说明有证据认为指标间并非完全不相关，不能推出``非常适合''，也不能确定只保留一个主成分。

KMO比较简单相关平方和与偏相关平方和的相对大小。按正文参考标准，0.42低于0.50，适合性不足；应检查所选指标的相关结构和测量内容。

保留成分数还需查看特征值、各成分及累计方差贡献率，并结合综合指标的专业含义。若第一成分保留信息不足或难以作合理解释，就不能只因球形检验显著而采用单一综合分。




\Answer{SUP-P3}{两科成绩的主成分与百分位评价}

\textbf{（1）表达式与方差。} 令 \(Z_1=(X_1-70)/10\)，\(Z_2=(X_2-80)/5\)，则 \[C_1=(Z_1+Z_2)/\sqrt2,\qquad C_2=(Z_1-Z_2)/\sqrt2.\] 对应特征值为1.8和0.2，贡献率为90\%和10\%，两成分不相关。

\textbf{（2）个体评价。} A的标准分为 \((2,1)\)，故 \(C_1=3/\sqrt2\approx2.1213\)，\(C_2=1/\sqrt2\approx0.7071\)。第一成分反映两科的综合相对水平；第二成分为正，表示A的语文相对优势高于数学相对优势，比较的是标准分。

\textbf{（3）百分位。} 第一成分的标准差是 \(\sqrt{1.8}\)，不是1。应先计算 \[F_1=\frac{C_1}{\sqrt{1.8}}\approx1.5811,\qquad \Phi(F_1)\approx0.9431.\] 即按该正态模型约处于第94.31百分位。这是模型给出的相对位置，不是A的百分制成绩，也不是直接数得的样本名次。

\textbf{（4）信息与条件。} 只保留第一成分会丢失占总标准化方差10\%的科目相对差异。PCA通常保证不相关；本题因另加联合正态条件，才可进一步说两个成分独立。部分资料所说的``独立''在一般资料中应按``不相关''理解。




\Answer{SUP-P4}{胎儿受精龄：还原主成分回归方程}

\textbf{（1）先合并标准分系数，再还原量纲。} \(\widehat Y=23.73+0.9484Z_1+1.0414Z_2+4.7536Z_3.\) 因此 \[\begin{aligned}
b_1&=0.9484/9.71\approx0.0976725,\\
b_2&=1.0414/6.86\approx0.1518076,\\
b_3&=4.7536/690.3\approx0.00688628,\\
b_0&=23.73-b_1(33.05)-b_2(23.26)-b_3(936.9)\approx10.51912.
\end{aligned}\] 故按题给系数得到 \[\boxed{\widehat Y\approx10.51912+0.0976725X_1+0.1518076X_2+0.00688628X_3.}\] 还原截距时必须扣除各均值项。原例另列简写回归方程；本题指定使用两位小数的成分系数，复算时不要求与末页简写结果逐位一致。

\textbf{（2）代入计算。} 标准化后得到 \(C_1\approx1.280864\)、\(C_2\approx0.098807\)，故 \(\widehat Y\approx29.0061\)周，约29.01周。用还原方程计算应得到相同结果。

\textbf{（3）方法含义。} 主成分彼此不相关，以保留的成分作自变量可避免原自变量的强线性相关造成的系数不稳定。舍弃成分同时舍弃了其中的信息，可能引入偏差；解释自变量方差少的方向也可能对结局有用，所以不能保证预测误差一定下降。
\EndExercises
